\documentclass[11pt]{article}

\usepackage{amsmath, amssymb}
\usepackage{geometry}
\usepackage{setspace}
\usepackage{titlesec}

\geometry{margin=1in}
\setstretch{1.1}

\title{Why Finite Learning Beings Need Constitutions:\\
Seven Observations on the Human Condition}

\author{}
\date{}

\begin{document}

\maketitle

\begin{abstract}
This paper offers seven observations any reflective human can make about
themselves: that the believer cannot be fully trusted, that revision is
possible, that the self depends on a substrate it did not make, that
its stakes were inherited rather than chosen, that learning passes
between generations, that present existence depends on others, and that
none of these conditions sustains itself without ongoing work. The
observations are presented plainly and made available for direct
verification. Taken together, they establish why constitutions, charters
of human rights, and the legal frameworks that surround them are not
optional artifacts but the structural form that collective written
guidance must take if it is to serve beings whose situation the
observations describe. Such documents are learning artifacts: they
accumulate constraints across generations in response to discovered
failures, distribute their coverage across many provisions to resist
single-point compromise, and remain open to revision as conditions
reveal what was not previously known. They are, in this sense, among
the most consequential things finite learning beings have built. The
paper is offered as a small contribution to recognising what such
documents are and why they are worth defending.
\end{abstract}

\section{Introduction}

This paper begins from a question. What can a reflective human say
about their own situation, using only what they can directly observe,
without appeal to authority, prior theory, or formal apparatus?

The question is asked deliberately. An understanding of the human
condition that depends on specialised knowledge cannot ground a claim
of generality. If only those who have read certain books, mastered
certain notations, or accepted certain authorities can verify the
understanding, then it is local rather than universal. It applies only
to the community already trained to accept it.

An understanding derivable from first-person observation is different.
Any reflective human capable of attending to their own experience can
verify or refute it. Such an understanding is then not the property of
any school or tradition. It is available to anyone willing to look
carefully at what they can directly observe.

The paper offers such an understanding, in the form of seven
observations developed in sequence. Each is presented plainly, made
available for direct verification, and developed with its immediate
implications.

The seven observations are not the whole of the paper. They are its
ground. From them, the paper turns to a question that the
observations make pressing: what kind of artifact is the appropriate
response to the situation the observations describe? The answer, we
argue, is a class of documents that has emerged across human history
in many forms — constitutions, charters of human rights, foundational
legal frameworks, and the bodies of law that surround them. These
documents are not contingent products of particular cultures. They
are the structural form that collective written guidance must take if
it is to serve beings whose situation the observations describe. The
paper closes with an argument that such documents are among the most
consequential things finite learning beings have built, and that
recognising what they are is part of what it takes to defend them.

The paper makes no claims to novelty, robustness, or completeness. It
claims only that the picture seemed worth articulating clearly enough
to be tested.

\section{On the Genesis of the Observations}

A reasonable reader will ask: if these observations are available to
any reflective being, why do they require articulation? Why are they
not common knowledge?

The answer is not that the observations require specialised knowledge
to be made. They do not. Anyone who attends honestly to their own
experience can verify each of the seven observations from what they
directly encounter. None of this requires a library.

The answer is rather that the observations are obscured. They are
obscured by inherited frameworks that present the world differently
from how it is. They are obscured by the experience of confidence,
which feels like access to truth even when it is not. They are
obscured by social arrangements that reward the performance of
certainty over the acknowledgement of doubt. They are obscured,
sometimes, by the language available to a thinker, which may not
contain words for what the thinker is encountering.

The author can speak to this from his own development. The first of
the observations, that the believing self cannot be trusted, was
reached in his early twenties as a personal recognition. It was not
derived from any text. It came from attending to his own experience
and noticing, with some surprise, that beliefs he had held with full
confidence had turned out to be wrong, and that the same feeling of
confidence accompanied beliefs he currently held. The recognition
began a quest, sustained over the following two decades, to articulate
what else could be observed about being a finite learning being.

The remaining observations were reached through this sustained
attention. Intellectual tools encountered along the way helped:
network theory provided a substrate-neutral language for thinking
about complexity; work in metacognition and the philosophy of mind
provided frameworks for recognising what was already encountered in
experience; complexity science and dynamical systems thinking
clarified the structural patterns being observed; lived experience,
including significant difficulty, supplied both the motivation to
keep looking and the texture of what was being observed.

The role of these intellectual tools was not to produce the
observations. The observations were already present in experience,
waiting to be articulated. The role of the tools was to provide
language for articulation and to help the observer work free of
inherited frames that had previously obscured what could be seen.
Different readers will have different frames to work free of, or
none at all. The point is not that any particular development is
required. The point is that the observations, once one has attended
to them, are not difficult to verify.

The paper's contribution is therefore one of articulation rather than
discovery. The observations have been available to any reflective
human for as long as there have been reflective humans. What the
paper offers is a compact set of statements, in plain language, that
the reader can check directly against their own experience. If the
reader finds the observations true, they were already true for the
reader before the reader read this paper. The paper has only made
them visible.

\section{Seven Observations}

\subsection{Observation One: Contradictory Belief}

A reflective being can notice, by attending to its own beliefs over
time, that it has held contradictory positions with equal confidence.
A belief held firmly yesterday may be held with equal firmness in its
opposite form today, or may be held simultaneously alongside its
contradiction without immediate awareness of the contradiction.
Anyone who has changed their mind about something important, or
caught themselves in inconsistency, has access to this observation.

The implication is that confidence is not a reliable indicator of
truth. The thinking subject cannot fully trust its own current
beliefs, because the same subjective experience of certainty has
previously accompanied beliefs the subject later rejected. This is
not a sceptical claim about all knowledge. It is an observation
about the structural unreliability of the believer.

\subsection{Observation Two: Learning}

The same reflective being can notice that it learns. Beliefs are
revised. Models are updated. Capacities are acquired that were not
present before. Children become adults. Novices become competent.
Mistakes are sometimes corrected. This is not a philosophical
position. It is a fact accessible to direct observation in oneself
and others.

Combined with the first observation, this yields a structural
description: a being whose current beliefs cannot be fully trusted,
but which is capable of revising those beliefs in response to
experience. The unreliability of belief is partially compensated by
the capacity for revision.

\subsection{Observation Three: Substrate}

The thinking self runs on something. The mind that observes its own
contradictions and learns from its own mistakes is not separable
from the body and brain that support it. Sleep, fatigue, injury,
intoxication, and aging all alter cognition in ways the thinking
subject cannot prevent through thought alone. Whatever the self is,
it is a configuration of an underlying substrate, and depends on
that substrate continuing to function.

This observation does not require any particular theory of
consciousness or mind-body relation. It requires only the
recognition that the conditions enabling thought are not themselves
products of thought. They are given, and they can be lost.

\subsection{Observation Four: Inherited Valence}

A reflective being can notice that its movements and decisions are
shaped by attractions and aversions that it did not choose. Warmth
is pleasant. Certain smells draw the body forward and others push it
away. Some situations are sought and others avoided. These
orientations are present before deliberation begins. The thinking
subject finds itself already attracted and already repelled, and its
decisions are made within the field these orientations create.

The observation is not that all behaviour is determined by such
orientations, but that they are present and were not authored by the
deliberating self. That warmth is valued is not a conclusion reached
by reasoning. It is given, by the substrate, before reasoning
begins.

The implication is that the substrate is not a neutral platform for
cognition. It supplies the orientations that make cognition
consequential. Without inherited attractions and aversions, there
would be no stakes to any outcome and no reason for the learning
process to continue rather than stop. The substrate does not merely
support the thinking self; it provides the directional pull that
gives thought its weight.

This observation also locates the thinking self within something
larger than itself. The orientations were not chosen; they were
inherited from the processes that produced the substrate.
Recognising this is recognising that the deliberating self sits
inside a body whose dispositions preceded it. What the self does
with these dispositions remains open. That they are there, and that
they shape what the self can do, is not.

\subsection{Observation Five: Generational Continuity}

A reflective being can observe that learning is not confined to a
single substrate. Knowledge passes between beings and across
generations. Parents teach children. Teachers instruct students.
Practices, languages, and skills survive the deaths of those who
held them, when the conditions for transmission are maintained.
Anyone who has been taught anything, or who has taught anything to
anyone else, has access to this observation.

The observation has two parts, both directly available to
experience.

The first is that substrates are mortal. Bodies age and fail.
Specific brains, with their specific configurations, do not persist
indefinitely. What was known by one substrate is lost when that
substrate fails, unless it has been transmitted to others.

The second is that learning nonetheless accumulates. Across
generations, what is known has grown. The understanding of the body
and brain that support thought took hundreds of years to develop and
required tools, theories, and methods that no single generation
could have produced alone. The progress of knowledge is equally
available as an observation: what is known today exceeds what any
individual could acquire from direct experience, because it has
been passed forward through substrates that did not survive but
whose learning did.

Together these establish that learning is generational. The
viability of any particular substrate is bounded; the viability of
learning across substrates is what allows knowledge to grow rather
than reset. This is not a claim that requires philosophical
defence. Anyone who reads this sentence is participating in such
transmission, using language they did not invent, in a script
developed by others, drawing on understanding accumulated across
generations who are no longer present.

\subsection{Observation Six: Interdependence}

A reflective being can observe that it does not produce most of
what it depends on. The food consumed today was grown by people
the eater has never met. The clothes worn were made by others. The
language used to think was developed by countless prior speakers
and continues to be shaped by current ones. The shelter, the tools,
the medicines, the infrastructure that allow the reflective being
to reflect at all are produced by a network of others whose work
the being cannot replace from its own resources.

This observation requires no theory and no instrument. It can be
made by anyone who looks at what they have eaten, worn, or used in
the past day, and asks where it came from.

Two implications follow.

The first is that the viability of the individual learner is not a
property of the individual alone. It depends on the continued
functioning of a network of others producing the conditions under
which the individual can continue to learn. If those conditions
fail, no degree of individual capacity will substitute for them. A
learner cut off from the network that feeds, clothes, and informs
them does not remain a learner for long.

The second is that solitary self-sufficiency, where it appears, is
itself a product of interdependence. The contemplative who believes
themselves independent is using language they did not invent,
sitting in shelter they did not build, having eaten food they did
not grow. The capacity to feel self-sufficient is itself produced
by the network whose existence the feeling denies.

\subsection{Observation Seven: Maintenance}

A reflective being can observe that the conditions enabling its
continuation do not sustain themselves. The house does not tidy
itself. The food does not arrive without being grown, harvested,
transported, and prepared. The body does not stay healthy without
sleep, eating, and care. The relationship does not survive without
attention. The skill does not remain sharp without practice.
Anyone who has run a household, raised a child, kept a friendship,
or maintained a body has direct access to this observation.

The observation has a sharp form: there is a minimum amount of
maintenance work required to keep any viable system viable, and
this minimum cannot be skipped indefinitely. Below it, the system
degrades. If the degradation continues, the system fails. The
failure ends what the maintenance was preserving.

This is most clearly experienced in caregiving relationships. A
parent discovers, often with some surprise, that there is no way
to sustain a household and a child without doing the chores. The
chores are not external impositions on the relationship; they are
constitutive of it. To stop doing them is to end the conditions
under which the relationship can continue. The same structure
appears in every domain of viability: bodies require sleep,
friendships require contact, institutions require administration,
languages require use, knowledge bases require updating.

Two implications follow.

The first is that viability is not a property a system has. It is
an ongoing accomplishment, produced by continuous low-level work
that often goes unnoticed precisely because it is successful. When
maintenance is performed, the system continues and the maintenance
becomes invisible. When maintenance fails, the collapse becomes
visible, but by then the conditions for repair may have been lost.

The second is that violating the maintenance floor ends the
relation it was sustaining. This is not a moral claim. It is a
structural one. A system that does not receive the minimum
required tending does not continue, regardless of the wishes or
intentions of those depending on it.

\section{Synthesis}

Read together, the seven observations yield a coherent picture of
what it is to be a finite learning being.

A being capable of reflection notices that it cannot fully trust
its own beliefs. It also notices that it can revise those beliefs
in response to experience. It depends, for its capacity to do
either, on a substrate it did not make and which it cannot fully
control. The orientations that give its decisions stakes were
inherited from that substrate before any decision was possible.
What it learns can outlive the particular substrate that learns,
if it is transmitted; what it currently depends on is produced by
a network of others performing work it cannot replace; and none of
these conditions sustains itself without ongoing tending.

This picture can be stated more compactly. A finite learning
being is a process running on a substrate it did not make, with
stakes it did not choose, dependent on others it cannot replace,
embedded in a generational chain it did not begin, and performing
maintenance it cannot skip.

Two consequences follow that prepare the next part of the paper.

The first is that no individual finite learning being can meet
this situation alone. The conditions identified by the
observations extend across timescales no individual life spans,
across networks no individual can fully see, and across futures
no individual can fully anticipate. Whatever response is
appropriate to the situation, it is not a response that can be
worked out and held within a single life.

The second is that the response, whatever it is, must itself
respect the conditions the observations describe. It must
acknowledge that its current articulation cannot be fully
trusted, that it must be open to revision, that it must be
transmissible across generations, that it must address the
dependencies among finite learning beings rather than only their
individual situations, and that it must include the work of its
own maintenance. A response that fails any of these criteria
will fail to serve the beings whose situation the observations
describe.

The next section turns to what such a response looks like.

\section{The Response: Constitutions and Charters as Living
Artifacts}

Across the long course of human history, finite learning beings
have built a particular kind of artifact in response to the
situation the observations describe. The artifact takes many
forms — constitutions, charters of human rights, foundational
legal documents, declarations of basic principles, the bodies of
law that surround and elaborate them. These forms differ in
their specifics, their origins, and the conditions under which
they were drafted. They share a structural function. They are
the form in which finite learning beings carry forward, across
generations, what they have learned about how to live together
under the conditions the observations describe.

We focus here on what such documents are, structurally, because
what they are matters as much as what they say.

\subsection{The Documents as Learning Artifacts}

A constitution or a charter of human rights is not a static text.
It is a learning artifact that grows over time. New constraints
are added when failures are detected that existing constraints
did not prevent. Old constraints are revised when their
thresholds turn out to be miscalibrated. The text as a whole
accumulates: each generation inherits a document that already
encodes the lessons of prior failures, and adds to it the
lessons learned in its own time. The document is, in this sense,
a written memory of what has been learned about how finite
learning beings can live together.

This dynamic structure is not an accident of legal practice. It
is the structural form that any collective written guidance must
take if it is to remain viable across the timescales the
observations require. A document that closes itself to revision
is a document that has stopped learning. A document that has
stopped learning cannot continue to serve finite learning beings
in conditions it did not anticipate. The seven observations apply
to the writers of such documents as much as to anyone else: the
current text cannot be fully trusted, the capacity for revision
must be preserved, and the maintenance work of keeping the
document current must be performed in every generation.

The text as it exists is the residue of this ongoing maintenance
work. Every constraint encoded in current foundational law was,
at some earlier moment, an observation by some generation that
something was going wrong, that a particular kind of failure was
recurring, and that a constraint was needed to address it. The
accumulation of constraints over time is the visible record of
collective learning about the conditions under which finite
learning beings can persist together.

\subsection{The Structural Form of Accumulation}

What such documents do, in structural terms, can be described as
follows. The space of possible failures of collective living is
large and not fully knowable in advance. Existing constraints
cover some of this space; others are not yet covered. When an
uncovered failure occurs, the failure itself constitutes evidence
that a new constraint is needed. The document grows by adding
constraints in response to such evidence, gradually extending
the coverage of the space of possible failures.

The growth is not arbitrary. New constraints are added where
failures have been detected, not where they might hypothetically
occur. This grounds the document in observed conditions rather
than in speculation. The accumulation is also gradual, with
constraints added one or a few at a time, allowing each addition
to be evaluated against subsequent experience. Constraints that
turn out to be poorly calibrated can be revised; constraints
that prove necessary become permanent features of the document.

The coverage is distributed across many provisions, no one of
which is sufficient on its own and no one of which is
indispensable. This distribution is itself a feature of the
document's viability. A document whose coverage rested on a
single provision would be fragile; the loss or compromise of
that provision would dissolve the whole. A document whose
coverage is distributed across many provisions, with
redundancies and overlapping protections, can absorb the loss
or revision of any particular provision without losing its
overall function. The robustness of such a document is a
property of its distributed structure, not of the strength of
any individual provision.

The same distributed structure provides resistance to capture.
A single bad actor, or a single misjudgment, cannot dissolve
the document, because no single point of failure controls it.
The document persists because it is structured to survive the
loss or compromise of parts of itself. This is what allows it
to operate under conditions of imperfect compliance, partial
enforcement, and ongoing contestation, none of which would be
survivable if its viability depended on perfect functioning of
any particular component.

\subsection{The Layered Structure of Foundational Documents}

Within the broader legal frameworks that finite learning beings
have built, foundational documents — constitutions, charters of
human rights, declarations of basic principles — play a
particular role. They articulate the deepest constraints, the
ones that the rest of the framework is required to respect.
They are typically harder to revise than ordinary law, and this
difficulty is not incidental: they carry the constraints that
have been identified as most fundamental to the conditions of
collective living, and the difficulty of revising them protects
against the short-term pressures that ordinary law must handle.

But foundational documents are not fixed documents. They are
slowly revising documents, whose revision rate is calibrated to
the timescale of the constraints they carry. Constraints at this
layer change rarely, but they do change, as the conditions they
address evolve and as new generations identify failures that the
existing text did not anticipate. Foundational documents are, in
this sense, the slowest learning components of the legal
frameworks they anchor: they revise less often than ordinary
law, but they revise.

The layered structure of legal frameworks reflects a structural
recognition: different constraints belong at different
timescales of revision. Constraints that should be stable across
many generations belong in foundational documents, where the
revision rate is slow and the threshold for change is high.
Constraints that must respond to changing conditions on shorter
timescales belong in ordinary law, where revision is more
frequent and the threshold is lower. Constraints that adapt to
particular cases and contexts belong in regulatory practice and
judicial interpretation, where revision is continuous. Each
layer operates at the timescale appropriate to the constraints
it carries. Together they perform, at the level of collective
living, what the seven observations describe at the level of
individual finite learning beings: the ongoing accumulation,
transmission, and revision of learning under conditions that no
individual or generation can fully control.

\subsection{Why Such Documents Are Worth Defending}

The seven observations make a case for why such documents are
not optional. Finite learning beings whose beliefs cannot be
trusted, whose substrates are mortal, whose stakes were
inherited rather than chosen, whose continuation depends on
others, and whose collective living requires maintenance no
individual can perform, need a written record of what has been
learned about their situation. They need it to outlive any
particular learner, to be testable across many lives, to address
structural features rather than particular cases, and to remain
open to revision as conditions reveal what was not previously
known.

This is what constitutions, charters of human rights, and the
legal frameworks they anchor are for. They are the form in which
finite learning beings carry forward, across generations, what
they have learned about how to live together. Each is an
imperfect, partial, sometimes badly compromised instance of this
form. Each was drafted under particular historical conditions
and bears the marks of those conditions. Each contains
constraints that are well-calibrated and constraints that have
turned out to be inadequate. None is finished. None is beyond
revision. None is the last word.

But each is also, in its own way, an extraordinary
accomplishment. The fact that finite learning beings have
managed to articulate, in written form, constraints on their
collective behaviour that have survived across generations, that
have been refined through repeated revision, that distribute
coverage across many provisions to resist single-point
compromise, and that remain open to the further revision the
observations require, is among the most consequential things
such beings have done. It is, in a precise sense, the human
condition learning about itself. The accumulated body of such
documents is the visible record of collective self-awareness:
of finite learning beings recognising their situation,
articulating constraints adequate to that situation, and passing
those constraints forward to the generations that will inherit
them.

To recognise these documents as such is to recognise what is at
stake in their preservation. They are not optional cultural
products that can be discarded when convenient. They are the
form in which the maintenance work of collective living, across
the timescales the observations require, is actually performed.
A generation that allows them to be dissolved, or stripped of
their distributed coverage, or closed to further revision, is a
generation that has stopped doing the maintenance work. The
observations are clear about what follows: the system that does
not receive the minimum required tending does not continue,
regardless of the wishes or intentions of those depending on
it.

To defend such documents is therefore not to defend any
particular text against any particular attack. It is to defend
the kind of artifact such documents are. It is to insist that
the work of collective written guidance, accumulated across
generations and revised as conditions reveal what was not
previously known, is the form in which finite learning beings
can live together under the conditions the observations
describe. The defence of such documents is, ultimately, the
defence of the conditions under which finite learning beings can
continue to be finite learning beings together.

\section{Limitations}

The picture offered here is small. It does not develop the
content of any particular foundational document. It does not
evaluate the constraints those documents carry. It does not
adjudicate between competing legal traditions. These tasks
belong to the substantive work of constitutional and legal
theory, and to the ongoing collective work of revising the
documents themselves.

What the paper offers is more modest: a recognition of what such
documents are, and why their existence is not optional for
beings whose situation the observations describe. The seven
observations ground the recognition. The structural analysis of
the documents draws out what follows from the observations
together with the recognition that any response to the
situation the observations describe must itself respect the
conditions the observations identify.

The picture is also not falsified merely by the existence of
counterexamples to any single observation. The claim is that the
seven observations, taken together, describe what it is to be a
finite learning being, and that this description grounds the
need for the kind of documents the paper has discussed. A reader
who finds one observation unconvincing in a particular case is
invited to consider whether the case is in fact an instance of
finite learning, or whether the observation requires refinement,
or whether the observation as stated does not generalise as
claimed. The picture is offered to support such examination,
not to forestall it.

\section{Conclusion}

This paper has offered seven observations about what it is to be
a finite learning being, derived from what any reflective human
can verify in their own experience. It has shown that these
observations together establish why a particular kind of artifact
— constitutions, charters of human rights, and the bodies of law
they anchor — is the structural form that collective written
guidance must take if it is to serve beings whose situation the
observations describe. Such documents are learning artifacts.
They accumulate constraints in response to discovered failures.
They distribute coverage across many provisions to resist
single-point compromise. They revise at timescales calibrated to
the depth of the constraints they carry. They remain open to
the further revision the observations require.

These documents are not optional. They are the form in which the
maintenance work of collective living is actually performed,
across the timescales no individual life spans. The accumulated
body of such documents is the visible record of finite learning
beings recognising their situation and articulating constraints
adequate to it. Each instance is imperfect; none is finished;
all are revisable. What matters is the kind of artifact such
documents are, and the recognition that defending them is part
of what it takes to continue living together under the
conditions the observations describe.

The paper offers no further argument. The observations are
available for verification. The structural analysis follows from
the observations together with the demand that the response
respect the conditions the observations identify. The reader is
invited to verify the observations from their own experience, to
test the structural analysis against their own knowledge of how
foundational legal documents have actually evolved, and to
consider what follows for the work of preserving and revising
such documents in their own time.

\end{document}